\subsection{Propósito e Arquitetura}
\label{sec:proposito_arquitetura}

Terminada a fase de aproximação de longa distância, inicia-se a aproximação de curta distância, onde o \gls{mr}, por meio de suas juntas revolutas, movimenta seus elos de modo que o efetuador se alinhe com a porta de acoplamento do alvo.
Durante essa fase, o comando das juntas é obtido a partir de um modelo dinâmico acoplado espaçonave--manipulador.

A lei de controle das juntas atua de forma integrada às leis de controle de atitude e do movimento relativo da espaçonave, assegurando que a orientação da espaçonave permaneça alinhada à do alvo, mesmo diante das reações geradas pela movimentação do manipulador.

Em ambiente de microgravidade, a conservação da quantidade de movimento angular total da espaçonave robótica. 
Assim, cada variação de configuração do \gls{mr} produz torques de reação que perturbam a atitude da espaçonave robótica, exigindo uma compensação ativa para manter a orientação desejada. 
Essa compensação é obtida por meio de uma lei de controle acoplado espaçonave robóticaada no modelo de Lagrange do conjunto espaçonave robótica--\gls{mr}, em que o tensor de inércia e os termos de Coriolis/centrífugos são utilizados para prever os torques de reação e alimentar as leis de controle de atitude.
O \gls{mr} adotado possui configuração 3R, composta por três juntas revolutas responsáveis pela movimentação e posicionamento da ferramenta, conforme exibido na Figura~\ref{fig:referenciais_MR}
